\section{考察}
\label{sec_discussion}
\subsection{精度を決めた要因}
研究の全体を通じて，精度を最も大きく変えたのはモデルの選択ではなく，ラベルの質とラベルの定義であった．
初期の比較では，同じ SigLIP2-so400m の埋め込みでも，メタデータの弱いラベルで学習すると正解率は約0.70，VLM のラベルで学習すると0.91であった．
VLM 自体も，プロンプトの定義を明確にするだけで正解率が約19ポイント向上したのに対し，8B から32B へのモデルの大型化では向上しなかった．
また，第2版から第3版への定義の改訂では，生徒の構成を変えずに test の3軸の平均が約0.07向上した．

ラベルの定義を改訂する際に有効であったのは，下流の処理から逆算する考え方と，判断が揃わない区別を無理に残さない判断である．
事前知識，合成画像，CG は，いずれも研究上の関心はあるが，判断者の知識や境界の解釈に依存し，教師と目視の判定が揃わなかった．
これらを廃止または統合した結果，残ったクラスはいずれも見た目で判定でき，かつ下流の処理と1対1に対応するものになった．
一方で，お題文と通常の境界は，形式の誤りの大半を占めるにもかかわらず，お題文の読み取りという処理に直結するため維持した．
精度のためにクラスを統合すると，マクロ F1 は機械的に上がるため，統合の判断は精度だけでなく用途に基づいて行う必要がある．

\subsection{VLM への指示の方法}
VLM への指示では，質問を分ける，対象を絞って別の短い質問で確かめるという方法が有効であり，説明を書き足す，思考の過程を書き出させるという方法は有効でなかった．
説明を書き足すと，Qwen3-VL-8B はかえって説明の対象のクラスを選びやすくなった．
また，見た目で区別しにくい概念を問い方の工夫で判定させる試みは，合成画像やお題文の体裁の質問のように，いずれも失敗した．
これらの結果は，VLM による分類のプロンプトは短く見た目に基づく基準で書き，判定の難しい部分は対象を絞った別の質問に切り出す方が安定することを示唆する．
ただし，これらの比較は Qwen3-VL-8B のみで行っており，他のモデルで同じ傾向が成り立つかは確かめていない．

\subsection{評価用データの重要性}
研究の過程では，評価用データで有効に見えた方法が，実データでは問題を起こす事例が2度あった．
1つ目は，媒体の定義を具体化したプロンプトであり，dev では媒体のマクロ F1 を 0.773 から 0.858 に上げたが，test では 0.859 から 0.760 に下げた．
差の大部分は各3枚以下の CG と絵画の正誤によるものであり，件数の少ない項目でプロンプトを調整すると偶然の正誤に引きずられる．
2つ目は，教師の構成の確認の質問であり，評価用データでは有効であったが，無作為の2万枚では典型的なお題文の多くを誤って外していた．
評価用データに典型的なテンプレートの画像が少なかったことが原因である．
これらの経験から，評価用データを dev と test に分けることに加え，少ない区分を多く含む評価用データと，実データの無作為抽出による確認を併用することが重要であるといえる．
層別抽出と抽出確率の逆数による重み付けは，少ないクラスの件数を確保しつつ実際の分布での精度を推定する方法として有用であった．

\subsection{限界}
本研究には次の限界がある．
第1に，全ての目視ラベルは1名の判断者によるものであり，判断者間の一致度を測っていない．
目視ラベルは，画像を表示して内容を確認できる AI エージェントが付与したものであり，人の判断者によるラベルとの一致は確かめていない．
特に，イラストとアニメやゲームの境界，テンプレートの文をお題文とするか，面積の約2割の領域に媒体を付けるかの判断には揺れが残る．
第2に，教師のプロンプトの多くは通常の評価用データの誤りを見て改訂したため，教師の精度は楽観的である．
生徒の構成の選択は dev で行い test で報告したが，プロンプトの選択の一部には test の結果も参照した段階がある．
第3に，絵画やゲームのような少ないクラスは，実データの目視ラベルでの件数が少なく，推定値の誤差が大きい．
第4に，評価用データの文字ありの見直しは，いずれかのモデルが文字ありと予測した画像に限ったため，全てのモデルが見落とした小さな文字が残っている可能性がある．
第5に，全データの集計は確認の質問の前の予測によるものであり，年ごとの精度の違いは評価していない．
